% Part: first-order-logic
% Chapter: syntax-and-semantics
% Section: substitution

\documentclass[../../../include/open-logic-section]{subfiles}

\begin{document}

\olfileid{fol}{syn}{sub}

\olsection{ځايناستي}

\begin{defn}[په ترم کښې ځايناستي]
$\Subst{s}{t}{x}$، يعنې په $s$ کښې د~$x$ د هر مورد پر ځاے د $t$
د \emph{ځايناستي} پايله، په بازګشتي ډول داسې تعريفوو:
\begin{enumerate}
\item \indcase{s}{c}{$\Subst{\indfrm}{t}{x}$ پخپله $s$ دے۔}

\item \indcase{s}{y}{که $y$ يو متغير او $y \not\ident x$ وي، نو
  $\Subst{\indfrm}{t}{x}$ هم پخپله~$s$ دے۔}

\item \indcase{s}{x}{$\Subst{\indfrm}{t}{x}$ همدا~$t$ دے۔}

\item\indcase{s}{\Atom{f}{t_1, \dots, t_n}}{$\Subst{\indfrmp}{t}{x}$
  همدا $\Atom{f}{\Subst{t_1}{t}{x}, \dots, \Subst{t_n}{t}{x}}$ دے۔}
\end{enumerate}
\end{defn}

\begin{defn}
يو ترم~$t$ په $!A$ کښې د $x$ دپاره هله \emph{!!{free for}} دے چې
په $!A$ کښې د~$x$ هېڅ ازاد مورد د داسې سور په ساحه کښې نۀ راځي چې
په~$t$ کښې کوم متغير تړي۔
\end{defn}

\begin{ex} ~
\begin{enumerate}
\item $\Obj v_8$ په $\lexists[\Obj
  v_3]\Atom{\Obj A^2_4}{\Obj v_3,\Obj v_1}$ کښې د $\Obj v_1$ دپاره
  ازاد دے۔

\item $\Obj f^2_1(\Obj v_1, \Obj v_2)$ په
  $\lforall[\Obj v_2]\Atom{\Obj A^2_4}{\Obj v_0,\Obj v_2}$ کښې د
  $\Obj v_0$ دپاره \emph{ازاد نۀ} دے۔
\end{enumerate}
\end{ex}

\begin{defn}[په !!a{formula} کښې ځايناستي]
که $!A$ !!a{formula}، $x$~!!a{variable}، او $t$~داسې ترم وي چې
په~$!A$ کښې د~$x$ دپاره !!{free for} وي، نو
$\Subst{!A}{t}{x}$ په~$!A$ کښې د~$x$ د ټولو ازادو موردونو پر ځاے
د~$t$ د ځايناستي پايله ده۔
\begin{enumerate}
\tagitem{prvFalse}{\indcase{!A}{\lfalse}{$\Subst{\indfrm}{t}{x}$
    همدا $\lfalse$ دے۔}}{}

\tagitem{prvTrue}{\indcase{!A}{\ltrue}{$\Subst{\indfrm}{t}{x}$
    همدا $\ltrue$ دے۔}}{}

\item \indcase{!A}{\Atom{P}{t_1,\dots,
    t_n}}{$\Subst{\indfrm}{t}{x}$ همدا
    $\Atom{P}{\Subst{t_1}{t}{x},
    \dots, \Subst{t_n}{t}{x}}$ دے۔}

\item \indcase{!A}{\eq[t_1][t_2]}{$\Subst{\indfrmp}{t}{x}$ همدا
  $\Subst{t_1}{t}{x} = \Subst{t_2}{t}{x}$ دے۔}

\tagitem{prvNot}{\indcase{!A}{\lnot !B}{$\Subst{\indfrmp}{t}{x}$
    همدا $\lnot \Subst{!B}{t}{x}$ دے۔}}{}

\tagitem{prvAnd}{\indcase{!A}{(!B \land
    !C)}{$\Subst{\indfrmp}{t}{x}$ همدا $(\Subst{!B}{t}{x} \land
    \Subst{!C}{t}{x})$ دے۔}}{}

\tagitem{prvOr}{\indcase{!A}{(!B \lor
    !C)}{$\Subst{\indfrmp}{t}{x}$ همدا $(\Subst{!B}{t}{x} \lor
    \Subst{!C}{t}{x})$ دے۔}}{}

\tagitem{prvIf}{\indcase{!A}{(!B \lif
    !C)}{$\Subst{\indfrmp}{t}{x}$ همدا $(\Subst{!B}{t}{x} \lif
    \Subst{!C}{t}{x})$ دے۔}}{}

\tagitem{prvIff}{\indcase{!A}{(!B \liff
    !C)}{$\Subst{\indfrmp}{t}{x}$ همدا $(\Subst{!B}{t}{x} \liff
    \Subst{!C}{t}{x})$ دے۔}}{}

\tagitem{prvAll}{
\indcase{!A}{\lforall[y][!B]}{که $y$ له $x$ څخه بېل متغير وي، نو
    $\Subst{\indfrmp}{t}{x}$ همدا
    $\lforall[y][\Subst{!B}{t}{x}]$ دے؛ که نۀ وي، نو
    $\Subst{\indfrmp}{t}{x}$ پخپله $\indfrm$ دے۔}}{}

\tagitem{prvEx}{
\indcase{!A}{\lexists[y][!B]}{که $y$ له $x$ څخه بېل متغير وي، نو
    $\Subst{\indfrmp}{t}{x}$ همدا
    $\lexists[y][\Subst{!B}{t}{x}]$ دے؛ که نۀ وي، نو
    $\Subst{\indfrmp}{t}{x}$ پخپله $\indfrm$ دے۔}}{}
\end{enumerate}
\end{defn}

\begin{explain}
پام وکړئ چې ځايناستي کېدے شي بې‌اغېزې وي: که $x$ په $!A$ کښې هېڅ
نۀ وي راغلے، نو $\Subst{!A}{t}{x}$ پخپله~$!A$ دے۔

دا محدوديت چې $t$ بايد په~$!A$ کښې د~$x$ دپاره !!{free for} وي، د
لاندې حالتونو د مخنيوي دپاره اړين دے۔ که $!A \ident
\lexists[y][x < y]$ او $t \ident y$ وي، نو $\Subst{!A}{t}{x}$ به
$\lexists[y][y < y]$ شي۔ په دې حالت کښې ازاد متغير $y$ د ځايناستي
پر مهال د سور $\lexists[y]$ له خوا «راګېرېږي»، او دا ناغوښتې پايله
ده۔ مثلاً غواړو چې هر کله $\lforall[x][!B]$ صدق کوي،
$\Subst{!B}{t}{x}$ هم صدق وکړي۔ خو
$\lforall[x][\lexists[y][x < y]]$ په پام کښې ونيسئ (دلته $!B$ همدا
$\lexists[y][x < y]$ دے)۔ دا د بېلګې په توګه د طبيعي عددونو په باب
رښتونې جمله ده: د هر عدد~$x$ دپاره تر هغه لوے يو عدد~$y$ شته۔ که
مو~$y$ د~$x$ د ممکنې ځايناستۍ په توګه روا ګڼلے واے، پايله به مو
$\Subst{!B}{y}{x} \ident \lexists[y][y < y]$ وه، چې دروغ ده۔ د دې
مخنيوے په دې کوو چې په~$t$ کښې هېڅ ازاد متغير بايد په~$!A$ کښې د
کوم سور له خوا تړلے نۀ شي۔
\end{explain}

د پېچلې نښې د مخنيوي دپاره ډېر ځله لاندې دود کاروو: که $!A$ داسې
!!a{formula} وي چې !!{variable}~$x$ پکښې ازاد راتلے شي، د دې د
ښودلو دپاره~$!A(x)$ هم ليکو۔ چې موخه مو کوم $!A$ او~$x$ دي څرګنده
وي، او $t$ يو ترم وي (چې فرض کېږي په $!A(x)$ کښې د $x$ دپاره ازاد
دے)، نو $!A(t)$ د $\Subst{!A}{t}{x}$ لنډون ليکو۔ مثلاً ويلے شو:
«$!A(t)$ د~$\lforall[x][!A(x)]$ يو مورد بولو۔» له دې څخه موخه دا
ده چې که $!A$ هر !!{formula}، $x$~!!a{variable}، او $t$~داسې ترم وي
چې په~$!A$ کښې د~$x$ دپاره ازاد وي، نو $\Subst{!A}{t}{x}$ د
$\lforall[x][!A]$ يو مورد دے۔

\end{document}
